\begin{ExerciseBlock}
\ExerciseHeading{07}{Python: Drehimpulsvektor}{angular_momentum_python_angular_momentum}
Schreiben Sie einen Python-Code, der die Aufgabe \textit{Drehimpulsvektor} aus diesem Übungsblatt löst. \\

Verwenden Sie die Vektorfunktionen aus dem Blatt zu Kräften als Ausgangspunkt.
\end{ExerciseBlock}

\begin{Solution}
\par\medskip\noindent\textbf{Musterlösung}\par\nopagebreak[4]
Berechnen Sie den Impuls aus $\vec p=m\vec v$ und bilden Sie das Kreuzprodukt $\vec L=\vec r\times\vec p$. Die Skalarprodukte prüfen die Orthogonalität. Der Winkel zwischen $\vec r$ und $\vec v$ folgt aus ihrem normierten Skalarprodukt.
\PythonCode{content/resources/chapters/angular_momentum/python/angular_momentum_vector.py}
\end{Solution}
